\section*{Safety and Scientific Validity}
\addcontentsline{toc}{section}{Safety and Scientific Validity}
\label{sec:front-matter-safety-and-scientific-validity}

The manual uses a bounded evidence ladder: build evidence, startup evidence, functional evidence, numerical evidence, scientific evidence, security evidence, and release readiness.

Figure~\ref{fig:manual-evidence-ladder} summarizes that ladder and the distinctions between the evidence levels.

% FIGURE-SPEC-BEGIN
% ID: fig:manual-evidence-ladder
% Source asset: figs/generated/manual-evidence-ladder
% Caption: Evidence ladder used by the manual.
% Accessibility: A hierarchical ladder that distinguishes build, startup, functional, numerical, scientific, security, and release evidence.
% FIGURE-SPEC-END
\begin{figure}[htbp]
\centering
\figureplaceholder{figs/generated/manual-evidence-ladder}
\caption{Evidence ladder used by the manual.}
\label{fig:manual-evidence-ladder}
\end{figure}

Figure~\ref{fig:manual-application-scope} summarizes the current application scope and the maturity level claimed for each family of tools.

% FIGURE-SPEC-BEGIN
% ID: fig:manual-application-scope
% Source asset: figs/generated/manual-application-scope
% Caption: Current application scope and maturity snapshot.
% Accessibility: A compact summary of which applications are validated for build or startup and which remain partially validated or planned.
% FIGURE-SPEC-END
\begin{figure}[htbp]
\centering
\figureplaceholder{figs/generated/manual-application-scope}
\caption{Current application scope and maturity snapshot.}
\label{fig:manual-application-scope}
\end{figure}

This book does not treat build success as scientific validation, and local-only or private-only services remain documented with that boundary intact.
